% !PRINTED 361
\weberpage{361}

\weq{8}{
\frac{m}{n} = (a_0, a_1, a_2 \ldots a_\nu)
= (a_0, a_1, a_2 \ldots a_\nu, 1).
}

I. \gesperrt{Es l\"asst sich also jeder rationale Bruch in einen Kettenbruch entwickeln, der nach Belieben eine gerade oder eine ungerade Anzahl von Theilnennern hat.}

Die in den Formeln (5) vorkommenden Zahlen
\[
\frac{m_1}{m_2}, \ \frac{m_2}{m_3}, \ \frac{m_3}{m_4} \ldots,
\]
nennen wir die \gesperrt{Schlusszahlen} des Kettenbruches, die Zahlen $a_0, a_1, a_2 \ldots$ sollen die \gesperrt{Theilnenner} heissen (obwohl $a_0$ nicht eigentlich als Nenner auftritt).

Ersetzen wir $m$ durch eine nach dem Modul $n$ congruente Zahl $m'$, so unterscheiden sich die beiden Br\"uche $m:n$ und $m':n$ nur um eine ganze Zahl; die Kettenbruchentwickelungen unterscheiden sich also nur in ihren ersten Gliedern $a_0$.

\begin{center}
\S.\,116.

\gesperrt{Kettenbruchentwickelung irrationaler Zahlen.}
\end{center}

Wenn $x$ eine reelle irrationale Zahl ist, so wird es immer eine und nur eine ganze Zahl $a_0$ geben, so dass $x$ zwischen $a_0$ und $a_0 + 1$ liegt. $a_0$ ist positiv, wenn $x$ gr\"osser als 1 ist, Null, wenn $x$ ein positiver echter Bruch ist, und negativ, wenn $x$ negativ ist. Setzen wir also
\[
x = a_0 + \frac{1}{x_1},
\]
so ist $x_1$ ein positiver unechter Bruch, also wenn $a_1$ die zun\"achst \gesperrt{unter} $x_1$ gelegene ganze Zahl ist, $a_1$ positiv. Wir setzen
\[
\begin{aligned}
x_1 &= a_1 + \frac{1}{x_2}, \\
x_2 &= a_2 + \frac{1}{x_3}, \\
&\ \cdots\cdots\cdots \\
x_{n-1} &= a_{n-1} + \frac{1}{x_n},
\end{aligned}
\]
%CONT>>
